\documentclass[12pt,a4paper]{article}
\usepackage[utf8]{vietnam}
\usepackage{amsmath,amssymb}
\usepackage{enumerate}
\usepackage{tabularx}
\usepackage{tikz}
\usepackage{xcolor}
\usepackage[top=0.45cm, bottom=0.45cm, left=0.8cm, right=0.8cm]{geometry}
\DeclareUnicodeCharacter{00B2}{\ensuremath{^2}}
\DeclareUnicodeCharacter{00B3}{\ensuremath{^3}}

\setlength{\parskip}{0pt}
\setlength{\parindent}{0pt}
\setlength{\fboxsep}{3.5pt}
\setlength{\fboxrule}{0.5pt}

% Dòng kẻ chấm mờ cự ly chuẩn viết tay tiểu học (18pt)
\newcommand{\dongke}[1]{%
  \begingroup%
  \setlength{\parskip}{0pt}%
  \baselineskip=18pt%
  \vspace{0.4mm}%
  \foreach \i in {1,...,#1}{%
    \noindent{\color{gray!55}\dotfill}\par%
  }%
  \vspace{0.2mm}%
  \endgroup%
}

% Khung bao quanh bài tập tự luận
\newcommand{\khungBaiTap}[1]{%
  \vspace{0.8mm}%
  \noindent\fbox{%
    \parbox{\dimexpr\linewidth-2\fboxsep-2\fboxrule\relax}{%
      \vspace{0.3mm}%
      #1%
      \vspace{0.3mm}%
    }%
  }%
  \vspace{0.5mm}%
}

% Ô bài tập đặt tính
\newcommand{\oDatTinh}[3][4]{%
  \vspace{0.3mm}%
  \noindent\textbf{#2} #3\par%
  \dongke{#1}%
}

% Câu trắc nghiệm dàn 4 đáp án trên 1 hàng ngang
\newcommand{\cauTN}[6]{%
  \noindent\textbf{Câu #1.} #2\par\vspace{0.1mm}%
  \noindent%
  \makebox[0.25\linewidth][l]{\textbf{A.} #3}%
  \makebox[0.25\linewidth][l]{\textbf{B.} #4}%
  \makebox[0.25\linewidth][l]{\textbf{C.} #5}%
  \makebox[0.25\linewidth][l]{\textbf{D.} #6}\par\vspace{0.6mm}%
}

% Câu trắc nghiệm dàn 2 đáp án / hàng
\newcommand{\cauTNhaiDong}[6]{%
  \noindent\textbf{Câu #1.} #2\par\vspace{0.1mm}%
  \noindent%
  \makebox[0.5\linewidth][l]{\textbf{A.} #3}%
  \makebox[0.5\linewidth][l]{\textbf{B.} #4}\par\vspace{0.1mm}%
  \noindent%
  \makebox[0.5\linewidth][l]{\textbf{C.} #5}%
  \makebox[0.5\linewidth][l]{\textbf{D.} #6}\par\vspace{0.6mm}%
}

\begin{document}
\pagestyle{empty}
\fontsize{10.2pt}{13.2pt}\selectfont
% =============================================================================
% ========================= MÃ ĐỀ: 501 ========================================
% =============================================================================

% -------------------- TRANG 1: MẶT TRƯỚC --------------------
\noindent
\begin{tabularx}{\linewidth}{|X|c|}
\hline
\textbf{PHIẾU BÀI TẬP CUỐI TUẦN -- TOÁN LỚP 5 (TUẦN 10)} & \textbf{MÃ ĐỀ: 501} \\
\textit{Trung tâm Toán Thầy Lĩnh -- Năm học: 2026 - 2027} & \textit{Thời gian: 45 phút} \\
\hline
\end{tabularx}

\vspace{1mm}

\noindent
\begin{tabularx}{\linewidth}{|X|p{3.5cm}|}
\hline
\rule[-2mm]{0pt}{7.5mm}\textbf{Họ và tên:} \dotfill & \textbf{Lớp:} \makebox[2cm]{\dotfill} \\
\hline
\end{tabularx}

\vspace{1mm}

\noindent
\begin{tabularx}{\linewidth}{|p{3cm}|X|}
\hline
\centering\textbf{ĐIỂM SỐ} & \centering\textbf{LỜI NHẬN XÉT CỦA THẦY CÔ} \tabularnewline
\hline
\rule{0pt}{1.1cm} & \tabularnewline
\hline
\end{tabularx}

\vspace{1mm}

\noindent\textbf{BẢNG ĐIỀN ĐÁP ÁN TRẮC NGHIỆM (Mỗi câu đúng 0,25 điểm):}
\vspace{0.6mm}

\begin{tabularx}{\linewidth}{|c|*{12}{>{\centering\arraybackslash}X|}}
\hline
\textbf{Câu} & \textbf{1} & \textbf{2} & \textbf{3} & \textbf{4} & \textbf{5} & \textbf{6} & \textbf{7} & \textbf{8} & \textbf{9} & \textbf{10} & \textbf{11} & \textbf{12} \\
\hline
\textbf{Đáp án} & & & & & & & & & & & & \\
\hline
\end{tabularx}

\vspace{1.2mm}

\noindent\textbf{\large PHẦN I. TRẮC NGHIỆM KHÁCH QUAN (3,0 điểm)}

\cauTN{1}{Kết quả của phép tính $25{,}36 + 18{,}45$ là:}{$43{,}81$}{$43{,}71$}{$44{,}81$}{$43{,}91$}
\cauTN{2}{Kết quả của phép tính $54{,}78 - 23{,}45$ là:}{$31{,}23$}{$31{,}33$}{$31{,}43$}{$30{,}33$}
\cauTN{3}{Kết quả của phép tính $72{,}4 - 35{,}28$ là:}{$37{,}22$}{$37{,}02$}{$37{,}12$}{$36{,}12$}
\cauTN{4}{Kết quả của phép tính $65 - 28{,}45$ là:}{$36{,}55$}{$37{,}55$}{$36{,}45$}{$37{,}45$}
\cauTN{5}{Biểu thức $a + b = b + a$ thể hiện tính chất nào:}{Tính chất kết hợp}{Tính chất giao hoán}{Tính chất phân phối}{Cộng với số 0}
\cauTN{6}{Giá trị của biểu thức $14{,}5 + 23{,}7 + 5{,}5$ là:}{$42{,}7$}{$43{,}5$}{$43{,}7$}{$44{,}7$}
\cauTN{7}{Tính nhanh tổng $3{,}7 + 4{,}25 + 6{,}3 + 5{,}75$ được kết quả là:}{$20$}{$21$}{$19{,}5$}{$20{,}5$}
\cauTN{8}{Tìm $x$, biết: $x + 12{,}5 = 35{,}8$. Giá trị của $x$ là:}{$22{,}3$}{$23{,}3$}{$24{,}3$}{$23{,}5$}
\cauTN{9}{Tìm $x$, biết: $48{,}6 - x = 19{,}4$. Giá trị của $x$ là:}{$28{,}2$}{$29{,}4$}{$29{,}2$}{$30{,}2$}
\cauTN{10}{Điền dấu thích hợp: $15{,}6 + 8{,}4 \dots 32{,}5 - 9{,}2$:}{$>$}{$<$}{$=$}{Không so sánh được}
\cauTNhaiDong{11}{Một xe tải chở hai bao gạo. Bao thứ nhất nặng $48{,}5\,\text{kg}$, bao thứ hai nặng hơn bao thứ nhất $6{,}25\,\text{kg}$. Cả hai bao nặng là:}{$102{,}25\,\text{kg}$}{$103{,}25\,\text{kg}$}{$104{,}25\,\text{kg}$}{$97{,}25\,\text{kg}$}
\cauTNhaiDong{12}{Một tam giác có độ dài ba cạnh lần lượt là $8{,}5\,\text{cm};\; 12{,}3\,\text{cm};\; 9{,}2\,\text{cm}$. Chu vi của tam giác đó là:}{$30\,\text{cm}$}{$29\,\text{cm}$}{$31\,\text{cm}$}{$30{,}5\,\text{cm}$}

\vspace{0.4mm}
\noindent\textbf{\large PHẦN II. TỰ LUẬN (7,0 điểm)}

\noindent\textbf{Bài 1 (2,0 điểm). Đặt tính rồi tính:}
\vspace{0.4mm}

\noindent
\begin{tabularx}{\linewidth}{|X|X|X|X|}
\hline
\oDatTinh[4]{a)}{$45{,}36 + 28{,}47$} & 
\oDatTinh[4]{b)}{$84{,}5 - 39{,}28$} & 
\oDatTinh[4]{c)}{$126{,}45 + 58{,}7$} & 
\oDatTinh[4]{d)}{$90 - 34{,}65$} \tabularnewline
\hline
\end{tabularx}

\newpage
% -------------------- TRANG 2: MẶT SAU --------------------
\noindent
\begin{tabularx}{\linewidth}{|X|c|}
\hline
\textbf{PHẦN TỰ LUẬN (TIẾP THEO) -- PHIẾU BÀI TẬP TUẦN 10} & \textbf{MÃ ĐỀ: 501} \\
\hline
\end{tabularx}

\vspace{0.8mm}

% BÀI 2: TÌM X (4 Ý)
\khungBaiTap{%
  \textbf{Bài 2 (1,5 điểm). Tìm $x$:}\par\vspace{0.6mm}
  \noindent
  \begin{minipage}[t]{0.485\linewidth}
    \textbf{a)} $x + 15{,}8 = 42{,}5$\par
    \dongke{3}
    \vspace{0.4mm}
    \textbf{c)} $64{,}2 - x = 27{,}5$\par
    \dongke{3}
  \end{minipage}\hfill
  \begin{minipage}[t]{0.485\linewidth}
    \textbf{b)} $x - 23{,}6 = 18{,}7$\par
    \dongke{3}
    \vspace{0.4mm}
    \textbf{d)} $x + 4{,}2 = 15{,}6 - 3{,}4$\par
    \dongke{3}
  \end{minipage}
}

\vspace{0.8mm}

% BÀI 3: ĐỔI ĐƠN VỊ (4 Ý - 1 CỘT TỪ TRÊN XUỐNG)
\khungBaiTap{%
  \textbf{Bài 3 (1,0 điểm). Viết số thích hợp vào chỗ chấm:}\par\vspace{0.8mm}
  \noindent
  a) $4\,\text{m}\,25\,\text{cm} = \dotfill\,\text{m}$\par\vspace{1.4mm}
  \noindent
  b) $7\,\text{km}\,350\,\text{m} = \dotfill\,\text{km}$\par\vspace{1.4mm}
  \noindent
  c) $5\,\text{kg}\,600\,\text{g} = \dotfill\,\text{kg}$\par\vspace{1.4mm}
  \noindent
  d) $3\,\text{tấn}\,45\,\text{kg} = \dotfill\,\text{tấn}$\par
}

\vspace{0.8mm}

% BÀI 4: BÀI TOÁN GIẢI (11 DÒNG KẺ)
\khungBaiTap{%
  \textbf{Bài 4 (1,5 điểm).} Một cửa hàng nhập về ba xe hàng. Xe thứ nhất chở 2{,}75 tấn, xe thứ hai chở nhiều hơn xe thứ nhất 0{,}6 tấn, xe thứ ba chở ít hơn tổng số hàng của hai xe đầu là 2{,}1 tấn. Hỏi cả ba xe chở bao nhiêu tấn hàng?
  \par\vspace{0.3mm}%
  \noindent\textit{Bài giải:}
  \dongke{11}
}

\vspace{0.8mm}

% BÀI 5: TOÁN NÂNG CAO (8 DÒNG KẺ)
\khungBaiTap{%
  \textbf{Bài 5 (1,0 điểm).} Tính bằng cách thuận tiện nhất:\par\vspace{0.6mm}\centerline{$A=3{,}7+4{,}25+6{,}3+5{,}75$}
  \par\vspace{0.3mm}%
  \noindent\textit{Bài làm:}
  \dongke{8}
}

\newpage

% =============================================================================
% ========================= MÃ ĐỀ: 502 ========================================
% =============================================================================

% -------------------- TRANG 1: MẶT TRƯỚC --------------------
\noindent
\begin{tabularx}{\linewidth}{|X|c|}
\hline
\textbf{PHIẾU BÀI TẬP CUỐI TUẦN -- TOÁN LỚP 5 (TUẦN 10)} & \textbf{MÃ ĐỀ: 502} \\
\textit{Trung tâm Toán Thầy Lĩnh -- Năm học: 2026 - 2027} & \textit{Thời gian: 45 phút} \\
\hline
\end{tabularx}

\vspace{1mm}

\noindent
\begin{tabularx}{\linewidth}{|X|p{3.5cm}|}
\hline
\rule[-2mm]{0pt}{7.5mm}\textbf{Họ và tên:} \dotfill & \textbf{Lớp:} \makebox[2cm]{\dotfill} \\
\hline
\end{tabularx}

\vspace{1mm}

\noindent
\begin{tabularx}{\linewidth}{|p{3cm}|X|}
\hline
\centering\textbf{ĐIỂM SỐ} & \centering\textbf{LỜI NHẬN XÉT CỦA THẦY CÔ} \tabularnewline
\hline
\rule{0pt}{1.1cm} & \tabularnewline
\hline
\end{tabularx}

\vspace{1mm}

\noindent\textbf{BẢNG ĐIỀN ĐÁP ÁN TRẮC NGHIỆM (Mỗi câu đúng 0,25 điểm):}
\vspace{0.6mm}

\begin{tabularx}{\linewidth}{|c|*{12}{>{\centering\arraybackslash}X|}}
\hline
\textbf{Câu} & \textbf{1} & \textbf{2} & \textbf{3} & \textbf{4} & \textbf{5} & \textbf{6} & \textbf{7} & \textbf{8} & \textbf{9} & \textbf{10} & \textbf{11} & \textbf{12} \\
\hline
\textbf{Đáp án} & & & & & & & & & & & & \\
\hline
\end{tabularx}

\vspace{1.2mm}

\noindent\textbf{\large PHẦN I. TRẮC NGHIỆM KHÁCH QUAN (3,0 điểm)}

\cauTN{1}{Kết quả của phép tính $36{,}48 + 27{,}35$ là:}{$63{,}83$}{$64{,}83$}{$63{,}73$}{$62{,}83$}
\cauTN{2}{Kết quả của phép tính $68{,}92 - 34{,}56$ là:}{$34{,}46$}{$34{,}36$}{$34{,}26$}{$33{,}36$}
\cauTN{3}{Kết quả của phép tính $83{,}5 - 46{,}37$ là:}{$37{,}23$}{$37{,}03$}{$37{,}13$}{$36{,}13$}
\cauTN{4}{Kết quả của phép tính $75 - 36{,}25$ là:}{$38{,}75$}{$39{,}75$}{$38{,}65$}{$39{,}65$}
\cauTN{5}{Biểu thức $(a + b) + c = a + (b + c)$ thể hiện tính chất nào:}{Tính chất giao hoán}{Tính chất kết hợp}{Tính chất phân phối}{Cộng với số 0}
\cauTN{6}{Giá trị của biểu thức $18{,}3 + 15{,}9 + 1{,}7$ là:}{$34{,}9$}{$35{,}8$}{$35{,}9$}{$36{,}9$}
\cauTN{7}{Tính nhanh tổng $4{,}6 + 5{,}35 + 5{,}4 + 4{,}65$ được kết quả là:}{$20$}{$21$}{$19{,}8$}{$20{,}5$}
\cauTN{8}{Tìm $x$, biết: $x + 18{,}4 = 42{,}6$. Giá trị của $x$ là:}{$23{,}2$}{$24{,}2$}{$25{,}2$}{$24{,}4$}
\cauTN{9}{Tìm $x$, biết: $56{,}8 - x = 28{,}3$. Giá trị của $x$ là:}{$27{,}5$}{$28{,}4$}{$28{,}5$}{$29{,}5$}
\cauTN{10}{Điền dấu thích hợp: $18{,}5 + 7{,}5 \dots 35{,}8 - 9{,}4$:}{$<$}{$>$}{$=$}{Không so sánh được}
\cauTNhaiDong{11}{Hai đoạn dây điện dài tất cả bao nhiêu nếu đoạn thứ nhất dài $24{,}8\,\text{m}$, đoạn thứ hai dài hơn đoạn thứ nhất $5{,}45\,\text{m}$:}{$54{,}05\,\text{m}$}{$55{,}05\,\text{m}$}{$56{,}05\,\text{m}$}{$50{,}25\,\text{m}$}
\cauTNhaiDong{12}{Một hình chữ nhật có nửa chu vi là $24{,}8\,\text{m}$. Chu vi của hình chữ nhật đó là:}{$49{,}6\,\text{m}$}{$48{,}6\,\text{m}$}{$50{,}6\,\text{m}$}{$47{,}6\,\text{m}$}

\vspace{0.4mm}
\noindent\textbf{\large PHẦN II. TỰ LUẬN (7,0 điểm)}

\noindent\textbf{Bài 1 (2,0 điểm). Đặt tính rồi tính:}
\vspace{0.4mm}

\noindent
\begin{tabularx}{\linewidth}{|X|X|X|X|}
\hline
\oDatTinh[4]{a)}{$56{,}48 + 35{,}26$} & 
\oDatTinh[4]{b)}{$92{,}4 - 45{,}38$} & 
\oDatTinh[4]{c)}{$135{,}62 + 47{,}8$} & 
\oDatTinh[4]{d)}{$80 - 28{,}75$} \tabularnewline
\hline
\end{tabularx}

\newpage
% -------------------- TRANG 2: MẶT SAU --------------------
\noindent
\begin{tabularx}{\linewidth}{|X|c|}
\hline
\textbf{PHẦN TỰ LUẬN (TIẾP THEO) -- PHIẾU BÀI TẬP TUẦN 10} & \textbf{MÃ ĐỀ: 502} \\
\hline
\end{tabularx}

\vspace{0.8mm}

% BÀI 2: TÌM X (4 Ý)
\khungBaiTap{%
  \textbf{Bài 2 (1,5 điểm). Tìm $x$:}\par\vspace{0.6mm}
  \noindent
  \begin{minipage}[t]{0.485\linewidth}
    \textbf{a)} $x + 18{,}5 = 45{,}2$\par
    \dongke{3}
    \vspace{0.4mm}
    \textbf{c)} $75{,}4 - x = 38{,}6$\par
    \dongke{3}
  \end{minipage}\hfill
  \begin{minipage}[t]{0.485\linewidth}
    \textbf{b)} $x - 27{,}4 = 21{,}8$\par
    \dongke{3}
    \vspace{0.4mm}
    \textbf{d)} $x + 5{,}6 = 18{,}4 - 4{,}2$\par
    \dongke{3}
  \end{minipage}
}

\vspace{0.8mm}

% BÀI 3: ĐỔI ĐƠN VỊ (4 Ý - 1 CỘT TỪ TRÊN XUỐNG)
\khungBaiTap{%
  \textbf{Bài 3 (1,0 điểm). Viết số thích hợp vào chỗ chấm:}\par\vspace{0.8mm}
  \noindent
  a) $5\,\text{m}\,48\,\text{cm} = \dotfill\,\text{m}$\par\vspace{1.4mm}
  \noindent
  b) $9\,\text{km}\,420\,\text{m} = \dotfill\,\text{km}$\par\vspace{1.4mm}
  \noindent
  c) $6\,\text{kg}\,350\,\text{g} = \dotfill\,\text{kg}$\par\vspace{1.4mm}
  \noindent
  d) $4\,\text{tấn}\,80\,\text{kg} = \dotfill\,\text{tấn}$\par
}

\vspace{0.8mm}

% BÀI 4: BÀI TOÁN GIẢI (11 DÒNG KẺ)
\khungBaiTap{%
  \textbf{Bài 4 (1,5 điểm).} Ba đội công nhân sửa đường. Đội thứ nhất sửa được 12{,}4 km, đội thứ hai sửa nhiều hơn đội thứ nhất 2{,}6 km, đội thứ ba sửa kém đội thứ hai 3{,}5 km. Hỏi cả ba đội sửa được bao nhiêu ki-lô-mét đường?
  \par\vspace{0.3mm}%
  \noindent\textit{Bài giải:}
  \dongke{11}
}

\vspace{0.8mm}

% BÀI 5: TOÁN NÂNG CAO (8 DÒNG KẺ)
\khungBaiTap{%
  \textbf{Bài 5 (1,0 điểm).} Tính bằng cách thuận tiện nhất:\par\vspace{0.6mm}\centerline{$B=12{,}5+8{,}6+7{,}5+1{,}4$}
  \par\vspace{0.3mm}%
  \noindent\textit{Bài làm:}
  \dongke{8}
}

\newpage

% =============================================================================
% ========================= MÃ ĐỀ: 503 ========================================
% =============================================================================

% -------------------- TRANG 1: MẶT TRƯỚC --------------------
\noindent
\begin{tabularx}{\linewidth}{|X|c|}
\hline
\textbf{PHIẾU BÀI TẬP CUỐI TUẦN -- TOÁN LỚP 5 (TUẦN 10)} & \textbf{MÃ ĐỀ: 503} \\
\textit{Trung tâm Toán Thầy Lĩnh -- Năm học: 2026 - 2027} & \textit{Thời gian: 45 phút} \\
\hline
\end{tabularx}

\vspace{1mm}

\noindent
\begin{tabularx}{\linewidth}{|X|p{3.5cm}|}
\hline
\rule[-2mm]{0pt}{7.5mm}\textbf{Họ và tên:} \dotfill & \textbf{Lớp:} \makebox[2cm]{\dotfill} \\
\hline
\end{tabularx}

\vspace{1mm}

\noindent
\begin{tabularx}{\linewidth}{|p{3cm}|X|}
\hline
\centering\textbf{ĐIỂM SỐ} & \centering\textbf{LỜI NHẬN XÉT CỦA THẦY CÔ} \tabularnewline
\hline
\rule{0pt}{1.1cm} & \tabularnewline
\hline
\end{tabularx}

\vspace{1mm}

\noindent\textbf{BẢNG ĐIỀN ĐÁP ÁN TRẮC NGHIỆM (Mỗi câu đúng 0,25 điểm):}
\vspace{0.6mm}

\begin{tabularx}{\linewidth}{|c|*{12}{>{\centering\arraybackslash}X|}}
\hline
\textbf{Câu} & \textbf{1} & \textbf{2} & \textbf{3} & \textbf{4} & \textbf{5} & \textbf{6} & \textbf{7} & \textbf{8} & \textbf{9} & \textbf{10} & \textbf{11} & \textbf{12} \\
\hline
\textbf{Đáp án} & & & & & & & & & & & & \\
\hline
\end{tabularx}

\vspace{1.2mm}

\noindent\textbf{\large PHẦN I. TRẮC NGHIỆM KHÁCH QUAN (3,0 điểm)}

\cauTN{1}{Kết quả của phép tính $42{,}57 + 31{,}68$ là:}{$74{,}25$}{$73{,}25$}{$74{,}15$}{$75{,}25$}
\cauTN{2}{Kết quả của phép tính $74{,}65 - 28{,}32$ là:}{$46{,}23$}{$46{,}33$}{$46{,}43$}{$45{,}33$}
\cauTN{3}{Kết quả của phép tính $65{,}3 - 27{,}46$ là:}{$37{,}94$}{$37{,}74$}{$37{,}84$}{$36{,}84$}
\cauTN{4}{Kết quả của phép tính $80 - 45{,}35$ là:}{$34{,}65$}{$35{,}65$}{$34{,}55$}{$35{,}55$}
\cauTN{5}{Khi đổi chỗ hai số hạng trong một tổng thì tổng đó:}{Tăng lên}{Không thay đổi}{Giảm đi}{Gấp đôi}
\cauTN{6}{Giá trị của biểu thức $23{,}4 + 16{,}8 + 6{,}6$ là:}{$45{,}8$}{$46{,}2$}{$46{,}8$}{$47{,}8$}
\cauTN{7}{Tính nhanh tổng $5{,}8 + 3{,}45 + 4{,}2 + 6{,}55$ được kết quả là:}{$20$}{$21$}{$19{,}5$}{$20{,}8$}
\cauTN{8}{Tìm $x$, biết: $x + 24{,}6 = 58{,}9$. Giá trị của $x$ là:}{$33{,}3$}{$34{,}3$}{$35{,}3$}{$34{,}5$}
\cauTN{9}{Tìm $x$, biết: $63{,}5 - x = 35{,}2$. Giá trị của $x$ là:}{$27{,}3$}{$28{,}5$}{$28{,}3$}{$29{,}3$}
\cauTN{10}{Điền dấu thích hợp: $24{,}8 + 6{,}2 \dots 45{,}6 - 14{,}5$:}{$<$}{$>$}{$=$}{Không so sánh được}
\cauTNhaiDong{11}{Một tấm vải dài $45\,\text{m}$. Buổi sáng bán được $15{,}5\,\text{m}$, buổi chiều bán được $18{,}75\,\text{m}$. Tấm vải còn lại là:}{$11{,}75\,\text{m}$}{$10{,}75\,\text{m}$}{$12{,}75\,\text{m}$}{$9{,}75\,\text{m}$}
\cauTNhaiDong{12}{Một tam giác có chu vi $36{,}8\,\text{cm}$, biết hai cạnh dài $11{,}5\,\text{cm}$ và $13{,}8\,\text{cm}$. Độ dài cạnh thứ ba là:}{$11{,}5\,\text{cm}$}{$12{,}5\,\text{cm}$}{$10{,}5\,\text{cm}$}{$11{,}8\,\text{cm}$}

\vspace{0.4mm}
\noindent\textbf{\large PHẦN II. TỰ LUẬN (7,0 điểm)}

\noindent\textbf{Bài 1 (2,0 điểm). Đặt tính rồi tính:}
\vspace{0.4mm}

\noindent
\begin{tabularx}{\linewidth}{|X|X|X|X|}
\hline
\oDatTinh[4]{a)}{$38{,}57 + 46{,}35$} & 
\oDatTinh[4]{b)}{$75{,}3 - 38{,}46$} & 
\oDatTinh[4]{c)}{$148{,}35 + 63{,}9$} & 
\oDatTinh[4]{d)}{$70 - 25{,}84$} \tabularnewline
\hline
\end{tabularx}

\newpage
% -------------------- TRANG 2: MẶT SAU --------------------
\noindent
\begin{tabularx}{\linewidth}{|X|c|}
\hline
\textbf{PHẦN TỰ LUẬN (TIẾP THEO) -- PHIẾU BÀI TẬP TUẦN 10} & \textbf{MÃ ĐỀ: 503} \\
\hline
\end{tabularx}

\vspace{0.8mm}

% BÀI 2: TÌM X (4 Ý)
\khungBaiTap{%
  \textbf{Bài 2 (1,5 điểm). Tìm $x$:}\par\vspace{0.6mm}
  \noindent
  \begin{minipage}[t]{0.485\linewidth}
    \textbf{a)} $x + 21{,}6 = 54{,}8$\par
    \dongke{3}
    \vspace{0.4mm}
    \textbf{c)} $82{,}6 - x = 45{,}8$\par
    \dongke{3}
  \end{minipage}\hfill
  \begin{minipage}[t]{0.485\linewidth}
    \textbf{b)} $x - 32{,}5 = 19{,}8$\par
    \dongke{3}
    \vspace{0.4mm}
    \textbf{d)} $x + 6{,}5 = 20{,}4 - 5{,}6$\par
    \dongke{3}
  \end{minipage}
}

\vspace{0.8mm}

% BÀI 3: ĐỔI ĐƠN VỊ (4 Ý - 1 CỘT TỪ TRÊN XUỐNG)
\khungBaiTap{%
  \textbf{Bài 3 (1,0 điểm). Viết số thích hợp vào chỗ chấm:}\par\vspace{0.8mm}
  \noindent
  a) $6\,\text{m}\,35\,\text{cm} = \dotfill\,\text{m}$\par\vspace{1.4mm}
  \noindent
  b) $5\,\text{km}\,680\,\text{m} = \dotfill\,\text{km}$\par\vspace{1.4mm}
  \noindent
  c) $7\,\text{kg}\,450\,\text{g} = \dotfill\,\text{kg}$\par\vspace{1.4mm}
  \noindent
  d) $5\,\text{tấn}\,120\,\text{kg} = \dotfill\,\text{tấn}$\par
}

\vspace{0.8mm}

% BÀI 4: BÀI TOÁN GIẢI (11 DÒNG KẺ)
\khungBaiTap{%
  \textbf{Bài 4 (1,5 điểm).} Một kho lương thực buổi sáng nhập 4{,}35 tấn gạo, buổi chiều nhập ít hơn buổi sáng 0{,}85 tấn. Sau đó kho xuất đi 5{,}2 tấn gạo. Hỏi kho còn lại bao nhiêu tấn gạo?
  \par\vspace{0.3mm}%
  \noindent\textit{Bài giải:}
  \dongke{11}
}

\vspace{0.8mm}

% BÀI 5: TOÁN NÂNG CAO (8 DÒNG KẺ)
\khungBaiTap{%
  \textbf{Bài 5 (1,0 điểm).} Tính bằng cách thuận tiện nhất:\par\vspace{0.6mm}\centerline{$C=7{,}25+3{,}8+2{,}75+6{,}2$}
  \par\vspace{0.3mm}%
  \noindent\textit{Bài làm:}
  \dongke{8}
}

\newpage

% =============================================================================
% ========================= MÃ ĐỀ: 504 ========================================
% =============================================================================

% -------------------- TRANG 1: MẶT TRƯỚC --------------------
\noindent
\begin{tabularx}{\linewidth}{|X|c|}
\hline
\textbf{PHIẾU BÀI TẬP CUỐI TUẦN -- TOÁN LỚP 5 (TUẦN 10)} & \textbf{MÃ ĐỀ: 504} \\
\textit{Trung tâm Toán Thầy Lĩnh -- Năm học: 2026 - 2027} & \textit{Thời gian: 45 phút} \\
\hline
\end{tabularx}

\vspace{1mm}

\noindent
\begin{tabularx}{\linewidth}{|X|p{3.5cm}|}
\hline
\rule[-2mm]{0pt}{7.5mm}\textbf{Họ và tên:} \dotfill & \textbf{Lớp:} \makebox[2cm]{\dotfill} \\
\hline
\end{tabularx}

\vspace{1mm}

\noindent
\begin{tabularx}{\linewidth}{|p{3cm}|X|}
\hline
\centering\textbf{ĐIỂM SỐ} & \centering\textbf{LỜI NHẬN XÉT CỦA THẦY CÔ} \tabularnewline
\hline
\rule{0pt}{1.1cm} & \tabularnewline
\hline
\end{tabularx}

\vspace{1mm}

\noindent\textbf{BẢNG ĐIỀN ĐÁP ÁN TRẮC NGHIỆM (Mỗi câu đúng 0,25 điểm):}
\vspace{0.6mm}

\begin{tabularx}{\linewidth}{|c|*{12}{>{\centering\arraybackslash}X|}}
\hline
\textbf{Câu} & \textbf{1} & \textbf{2} & \textbf{3} & \textbf{4} & \textbf{5} & \textbf{6} & \textbf{7} & \textbf{8} & \textbf{9} & \textbf{10} & \textbf{11} & \textbf{12} \\
\hline
\textbf{Đáp án} & & & & & & & & & & & & \\
\hline
\end{tabularx}

\vspace{1.2mm}

\noindent\textbf{\large PHẦN I. TRẮC NGHIỆM KHÁCH QUAN (3,0 điểm)}

\cauTN{1}{Kết quả của phép tính $53{,}64 + 28{,}79$ là:}{$82{,}43$}{$81{,}43$}{$82{,}33$}{$83{,}43$}
\cauTN{2}{Kết quả của phép tính $82{,}54 - 36{,}21$ là:}{$46{,}23$}{$46{,}33$}{$46{,}43$}{$45{,}33$}
\cauTN{3}{Kết quả của phép tính $91{,}2 - 45{,}68$ là:}{$45{,}62$}{$45{,}42$}{$45{,}52$}{$44{,}52$}
\cauTN{4}{Kết quả của phép tính $100 - 54{,}35$ là:}{$45{,}65$}{$46{,}65$}{$45{,}55$}{$46{,}55$}
\cauTN{5}{Muốn cộng hai số thập phân, ta viết số hạng này dưới số hạng kia sao cho:}{Các chữ số thẳng hàng nhau}{Các chữ số ở cùng một hàng đặt thẳng cột}{Chữ số đầu tiên thẳng hàng}{Tùy ý}
\cauTN{6}{Giá trị của biểu thức $27{,}6 + 18{,}5 + 2{,}4$ là:}{$47{,}5$}{$48{,}1$}{$48{,}5$}{$49{,}5$}
\cauTN{7}{Tính nhanh tổng $6{,}2 + 4{,}75 + 3{,}8 + 5{,}25$ được kết quả là:}{$20$}{$21$}{$19{,}6$}{$20{,}5$}
\cauTN{8}{Tìm $x$, biết: $x + 28{,}5 = 64{,}8$. Giá trị của $x$ là:}{$35{,}3$}{$36{,}3$}{$37{,}3$}{$36{,}5$}
\cauTN{9}{Tìm $x$, biết: $72{,}4 - x = 41{,}6$. Giá trị của $x$ là:}{$29{,}8$}{$31{,}2$}{$30{,}8$}{$30{,}2$}
\cauTN{10}{Điền dấu thích hợp: $28{,}6 + 11{,}4 \dots 52{,}4 - 12{,}5$:}{$>$}{$<$}{$=$}{Không so sánh được}
\cauTNhaiDong{11}{Một bể đang có $120{,}5\,\text{lít}$ nước. Người ta đổ thêm vào hai thùng, thùng thứ nhất $25{,}5\,\text{lít}$, thùng thứ hai $28{,}75\,\text{lít}$. Trong bể có tất cả số lít nước là:}{$173{,}75\,\text{lít}$}{$174{,}75\,\text{lít}$}{$175{,}75\,\text{lít}$}{$172{,}25\,\text{lít}$}
\cauTNhaiDong{12}{Một hình tam giác có độ dài đáy $15{,}4\,\text{m}$, chiều cao kém đáy $3{,}6\,\text{m}$. Chiều cao của hình tam giác đó là:}{$11{,}8\,\text{m}$}{$12{,}8\,\text{m}$}{$11{,}2\,\text{m}$}{$12{,}2\,\text{m}$}

\vspace{0.4mm}
\noindent\textbf{\large PHẦN II. TỰ LUẬN (7,0 điểm)}

\noindent\textbf{Bài 1 (2,0 điểm). Đặt tính rồi tính:}
\vspace{0.4mm}

\noindent
\begin{tabularx}{\linewidth}{|X|X|X|X|}
\hline
\oDatTinh[4]{a)}{$67{,}29 + 25{,}84$} & 
\oDatTinh[4]{b)}{$86{,}2 - 49{,}35$} & 
\oDatTinh[4]{c)}{$154{,}28 + 39{,}8$} & 
\oDatTinh[4]{d)}{$100 - 46{,}35$} \tabularnewline
\hline
\end{tabularx}

\newpage
% -------------------- TRANG 2: MẶT SAU --------------------
\noindent
\begin{tabularx}{\linewidth}{|X|c|}
\hline
\textbf{PHẦN TỰ LUẬN (TIẾP THEO) -- PHIẾU BÀI TẬP TUẦN 10} & \textbf{MÃ ĐỀ: 504} \\
\hline
\end{tabularx}

\vspace{0.8mm}

% BÀI 2: TÌM X (4 Ý)
\khungBaiTap{%
  \textbf{Bài 2 (1,5 điểm). Tìm $x$:}\par\vspace{0.6mm}
  \noindent
  \begin{minipage}[t]{0.485\linewidth}
    \textbf{a)} $x + 25{,}7 = 63{,}4$\par
    \dongke{3}
    \vspace{0.4mm}
    \textbf{c)} $92{,}4 - x = 48{,}7$\par
    \dongke{3}
  \end{minipage}\hfill
  \begin{minipage}[t]{0.485\linewidth}
    \textbf{b)} $x - 38{,}6 = 24{,}9$\par
    \dongke{3}
    \vspace{0.4mm}
    \textbf{d)} $x + 7{,}8 = 24{,}5 - 6{,}9$\par
    \dongke{3}
  \end{minipage}
}

\vspace{0.8mm}

% BÀI 3: ĐỔI ĐƠN VỊ (4 Ý - 1 CỘT TỪ TRÊN XUỐNG)
\khungBaiTap{%
  \textbf{Bài 3 (1,0 điểm). Viết số thích hợp vào chỗ chấm:}\par\vspace{0.8mm}
  \noindent
  a) $7\,\text{m}\,65\,\text{cm} = \dotfill\,\text{m}$\par\vspace{1.4mm}
  \noindent
  b) $8\,\text{km}\,250\,\text{m} = \dotfill\,\text{km}$\par\vspace{1.4mm}
  \noindent
  c) $4\,\text{kg}\,800\,\text{g} = \dotfill\,\text{kg}$\par\vspace{1.4mm}
  \noindent
  d) $6\,\text{tấn}\,35\,\text{kg} = \dotfill\,\text{tấn}$\par
}

\vspace{0.8mm}

% BÀI 4: BÀI TOÁN GIẢI (11 DÒNG KẺ)
\khungBaiTap{%
  \textbf{Bài 4 (1,5 điểm).} Quãng đường từ nhà An đến trường dài 2{,}6 km. An đi bộ được 0{,}75 km rồi được bố chở tiếp thêm 1{,}5 km. Đoạn đường còn lại An tự đi. Hỏi An còn phải đi bao nhiêu ki-lô-mét nữa?
  \par\vspace{0.3mm}%
  \noindent\textit{Bài giải:}
  \dongke{11}
}

\vspace{0.8mm}

% BÀI 5: TOÁN NÂNG CAO (8 DÒNG KẺ)
\khungBaiTap{%
  \textbf{Bài 5 (1,0 điểm).} Tính bằng cách thuận tiện nhất:\par\vspace{0.6mm}\centerline{$D=15{,}4-2{,}7-7{,}3$}
  \par\vspace{0.3mm}%
  \noindent\textit{Bài làm:}
  \dongke{8}
}

\newpage

% =============================================================================
% ========================= MÃ ĐỀ: 505 ========================================
% =============================================================================

% -------------------- TRANG 1: MẶT TRƯỚC --------------------
\noindent
\begin{tabularx}{\linewidth}{|X|c|}
\hline
\textbf{PHIẾU BÀI TẬP CUỐI TUẦN -- TOÁN LỚP 5 (TUẦN 10)} & \textbf{MÃ ĐỀ: 505} \\
\textit{Trung tâm Toán Thầy Lĩnh -- Năm học: 2026 - 2027} & \textit{Thời gian: 45 phút} \\
\hline
\end{tabularx}

\vspace{1mm}

\noindent
\begin{tabularx}{\linewidth}{|X|p{3.5cm}|}
\hline
\rule[-2mm]{0pt}{7.5mm}\textbf{Họ và tên:} \dotfill & \textbf{Lớp:} \makebox[2cm]{\dotfill} \\
\hline
\end{tabularx}

\vspace{1mm}

\noindent
\begin{tabularx}{\linewidth}{|p{3cm}|X|}
\hline
\centering\textbf{ĐIỂM SỐ} & \centering\textbf{LỜI NHẬN XÉT CỦA THẦY CÔ} \tabularnewline
\hline
\rule{0pt}{1.1cm} & \tabularnewline
\hline
\end{tabularx}

\vspace{1mm}

\noindent\textbf{BẢNG ĐIỀN ĐÁP ÁN TRẮC NGHIỆM (Mỗi câu đúng 0,25 điểm):}
\vspace{0.6mm}

\begin{tabularx}{\linewidth}{|c|*{12}{>{\centering\arraybackslash}X|}}
\hline
\textbf{Câu} & \textbf{1} & \textbf{2} & \textbf{3} & \textbf{4} & \textbf{5} & \textbf{6} & \textbf{7} & \textbf{8} & \textbf{9} & \textbf{10} & \textbf{11} & \textbf{12} \\
\hline
\textbf{Đáp án} & & & & & & & & & & & & \\
\hline
\end{tabularx}

\vspace{1.2mm}

\noindent\textbf{\large PHẦN I. TRẮC NGHIỆM KHÁCH QUAN (3,0 điểm)}

\cauTN{1}{Kết quả của phép tính $48{,}75 + 34{,}68$ là:}{$83{,}43$}{$82{,}43$}{$83{,}33$}{$84{,}43$}
\cauTN{2}{Kết quả của phép tính $76{,}85 - 32{,}48$ là:}{$44{,}27$}{$44{,}37$}{$44{,}47$}{$43{,}37$}
\cauTN{3}{Kết quả của phép tính $84{,}5 - 38{,}62$ là:}{$45{,}98$}{$45{,}78$}{$45{,}88$}{$44{,}88$}
\cauTN{4}{Kết quả của phép tính $85 - 39{,}48$ là:}{$45{,}52$}{$46{,}52$}{$45{,}42$}{$46{,}42$}
\cauTN{5}{Khi cộng nhiều số thập phân ta có thể:}{Chỉ cộng từ trái sang phải}{Đổi chỗ và nhóm các số hạng một cách thuận tiện}{Không được đổi chỗ}{Luôn phải đặt tính theo cột}
\cauTN{6}{Giá trị của biểu thức $32{,}5 + 14{,}8 + 7{,}5$ là:}{$53{,}8$}{$54{,}2$}{$54{,}8$}{$55{,}8$}
\cauTN{7}{Tính nhanh tổng $7{,}4 + 3{,}65 + 2{,}6 + 6{,}35$ được kết quả là:}{$20$}{$21$}{$19{,}5$}{$20{,}4$}
\cauTN{8}{Tìm $x$, biết: $x + 31{,}4 = 75{,}6$. Giá trị của $x$ là:}{$43{,}2$}{$44{,}2$}{$45{,}2$}{$44{,}5$}
\cauTN{9}{Tìm $x$, biết: $81{,}5 - x = 46{,}8$. Giá trị của $x$ là:}{$33{,}7$}{$35{,}2$}{$34{,}7$}{$34{,}2$}
\cauTN{10}{Điền dấu thích hợp: $34{,}5 + 12{,}5 \dots 65{,}2 - 17{,}8$:}{$<$}{$>$}{$=$}{Không so sánh được}
\cauTNhaiDong{11}{Một cửa hàng có $350\,\text{kg}$ gạo. Lần thứ nhất bán $85{,}5\,\text{kg}$, lần thứ hai bán $112{,}75\,\text{kg}$. Sau hai lần bán cửa hàng còn lại số gạo là:}{$150{,}75\,\text{kg}$}{$151{,}75\,\text{kg}$}{$152{,}75\,\text{kg}$}{$149{,}75\,\text{kg}$}
\cauTNhaiDong{12}{Một mảnh đất hình chữ nhật có chiều dài $28{,}5\,\text{m}$, chiều rộng ngắn hơn chiều dài $7{,}8\,\text{m}$. Chiều rộng mảnh đất là:}{$20{,}7\,\text{m}$}{$21{,}7\,\text{m}$}{$20{,}2\,\text{m}$}{$21{,}2\,\text{m}$}

\vspace{0.4mm}
\noindent\textbf{\large PHẦN II. TỰ LUẬN (7,0 điểm)}

\noindent\textbf{Bài 1 (2,0 điểm). Đặt tính rồi tính:}
\vspace{0.4mm}

\noindent
\begin{tabularx}{\linewidth}{|X|X|X|X|}
\hline
\oDatTinh[4]{a)}{$49{,}65 + 38{,}58$} & 
\oDatTinh[4]{b)}{$95{,}4 - 58{,}27$} & 
\oDatTinh[4]{c)}{$118{,}74 + 75{,}6$} & 
\oDatTinh[4]{d)}{$85 - 39{,}48$} \tabularnewline
\hline
\end{tabularx}

\newpage
% -------------------- TRANG 2: MẶT SAU --------------------
\noindent
\begin{tabularx}{\linewidth}{|X|c|}
\hline
\textbf{PHẦN TỰ LUẬN (TIẾP THEO) -- PHIẾU BÀI TẬP TUẦN 10} & \textbf{MÃ ĐỀ: 505} \\
\hline
\end{tabularx}

\vspace{0.8mm}

% BÀI 2: TÌM X (4 Ý)
\khungBaiTap{%
  \textbf{Bài 2 (1,5 điểm). Tìm $x$:}\par\vspace{0.6mm}
  \noindent
  \begin{minipage}[t]{0.485\linewidth}
    \textbf{a)} $x + 28{,}6 = 72{,}4$\par
    \dongke{3}
    \vspace{0.4mm}
    \textbf{c)} $86{,}5 - x = 39{,}7$\par
    \dongke{3}
  \end{minipage}\hfill
  \begin{minipage}[t]{0.485\linewidth}
    \textbf{b)} $x - 34{,}8 = 27{,}5$\par
    \dongke{3}
    \vspace{0.4mm}
    \textbf{d)} $x + 8{,}4 = 25{,}6 - 7{,}5$\par
    \dongke{3}
  \end{minipage}
}

\vspace{0.8mm}

% BÀI 3: ĐỔI ĐƠN VỊ (4 Ý - 1 CỘT TỪ TRÊN XUỐNG)
\khungBaiTap{%
  \textbf{Bài 3 (1,0 điểm). Viết số thích hợp vào chỗ chấm:}\par\vspace{0.8mm}
  \noindent
  a) $8\,\text{m}\,45\,\text{cm} = \dotfill\,\text{m}$\par\vspace{1.4mm}
  \noindent
  b) $6\,\text{km}\,150\,\text{m} = \dotfill\,\text{km}$\par\vspace{1.4mm}
  \noindent
  c) $9\,\text{kg}\,250\,\text{g} = \dotfill\,\text{kg}$\par\vspace{1.4mm}
  \noindent
  d) $7\,\text{tấn}\,50\,\text{kg} = \dotfill\,\text{tấn}$\par
}

\vspace{0.8mm}

% BÀI 4: BÀI TOÁN GIẢI (11 DÒNG KẺ)
\khungBaiTap{%
  \textbf{Bài 4 (1,5 điểm).} Ba bao gạo có khối lượng lần lượt là 12{,}5 kg, 15{,}75 kg và 10{,}75 kg. Người ta lấy ra từ bao thứ nhất 2{,}5 kg và thêm vào bao thứ ba 3{,}25 kg. Hỏi lúc này cả ba bao có bao nhiêu ki-lô-gam gạo?
  \par\vspace{0.3mm}%
  \noindent\textit{Bài giải:}
  \dongke{11}
}

\vspace{0.8mm}

% BÀI 5: TOÁN NÂNG CAO (8 DÒNG KẺ)
\khungBaiTap{%
  \textbf{Bài 5 (1,0 điểm).} Tính bằng cách thuận tiện nhất:\par\vspace{0.6mm}\centerline{$E=9{,}25+6{,}4+0{,}75+3{,}6$}
  \par\vspace{0.3mm}%
  \noindent\textit{Bài làm:}
  \dongke{8}
}

\newpage

% =============================================================================
% ========================= MÃ ĐỀ: 506 ========================================
% =============================================================================

% -------------------- TRANG 1: MẶT TRƯỚC --------------------
\noindent
\begin{tabularx}{\linewidth}{|X|c|}
\hline
\textbf{PHIẾU BÀI TẬP CUỐI TUẦN -- TOÁN LỚP 5 (TUẦN 10)} & \textbf{MÃ ĐỀ: 506} \\
\textit{Trung tâm Toán Thầy Lĩnh -- Năm học: 2026 - 2027} & \textit{Thời gian: 45 phút} \\
\hline
\end{tabularx}

\vspace{1mm}

\noindent
\begin{tabularx}{\linewidth}{|X|p{3.5cm}|}
\hline
\rule[-2mm]{0pt}{7.5mm}\textbf{Họ và tên:} \dotfill & \textbf{Lớp:} \makebox[2cm]{\dotfill} \\
\hline
\end{tabularx}

\vspace{1mm}

\noindent
\begin{tabularx}{\linewidth}{|p{3cm}|X|}
\hline
\centering\textbf{ĐIỂM SỐ} & \centering\textbf{LỜI NHẬN XÉT CỦA THẦY CÔ} \tabularnewline
\hline
\rule{0pt}{1.1cm} & \tabularnewline
\hline
\end{tabularx}

\vspace{1mm}

\noindent\textbf{BẢNG ĐIỀN ĐÁP ÁN TRẮC NGHIỆM (Mỗi câu đúng 0,25 điểm):}
\vspace{0.6mm}

\begin{tabularx}{\linewidth}{|c|*{12}{>{\centering\arraybackslash}X|}}
\hline
\textbf{Câu} & \textbf{1} & \textbf{2} & \textbf{3} & \textbf{4} & \textbf{5} & \textbf{6} & \textbf{7} & \textbf{8} & \textbf{9} & \textbf{10} & \textbf{11} & \textbf{12} \\
\hline
\textbf{Đáp án} & & & & & & & & & & & & \\
\hline
\end{tabularx}

\vspace{1.2mm}

\noindent\textbf{\large PHẦN I. TRẮC NGHIỆM KHÁCH QUAN (3,0 điểm)}

\cauTN{1}{Kết quả của phép tính $62{,}45 + 28{,}76$ là:}{$91{,}21$}{$90{,}21$}{$91{,}11$}{$92{,}21$}
\cauTN{2}{Kết quả của phép tính $85{,}74 - 42{,}38$ là:}{$43{,}26$}{$43{,}36$}{$43{,}46$}{$42{,}36$}
\cauTN{3}{Kết quả của phép tính $78{,}5 - 36{,}48$ là:}{$42{,}12$}{$41{,}92$}{$42{,}02$}{$41{,}02$}
\cauTN{4}{Kết quả của phép tính $95 - 42{,}65$ là:}{$52{,}35$}{$53{,}35$}{$52{,}25$}{$53{,}25$}
\cauTN{5}{Giá trị của $x - 0$ bằng:}{$0$}{$x$}{$2x$}{$x + 1$}
\cauTN{6}{Giá trị của biểu thức $36{,}8 + 21{,}4 + 3{,}2$ là:}{$60{,}4$}{$61{,}2$}{$61{,}4$}{$62{,}4$}
\cauTN{7}{Tính nhanh tổng $8{,}3 + 2{,}45 + 1{,}7 + 7{,}55$ được kết quả là:}{$20$}{$21$}{$19{,}8$}{$20{,}5$}
\cauTN{8}{Tìm $x$, biết: $x + 35{,}8 = 82{,}4$. Giá trị của $x$ là:}{$45{,}6$}{$46{,}6$}{$47{,}6$}{$46{,}8$}
\cauTN{9}{Tìm $x$, biết: $94{,}6 - x = 52{,}3$. Giá trị của $x$ là:}{$41{,}3$}{$43{,}3$}{$42{,}3$}{$42{,}5$}
\cauTN{10}{Điền dấu thích hợp: $42{,}8 + 15{,}2 \dots 75{,}6 - 16{,}8$:}{$<$}{$>$}{$=$}{Không so sánh được}
\cauTNhaiDong{11}{Ba bạn An, Bình, Cường cân nặng tất cả bao nhiêu nếu An nặng $38{,}5\,\text{kg}$, Bình nặng hơn An $2{,}3\,\text{kg}$, Cường nặng $41{,}2\,\text{kg}$:}{$119{,}5\,\text{kg}$}{$120{,}5\,\text{kg}$}{$121{,}5\,\text{kg}$}{$118{,}5\,\text{kg}$}
\cauTNhaiDong{12}{Một sợi dây dài $18{,}4\,\text{m}$. Người ta cắt đi một đoạn dài $5{,}65\,\text{m}$. Đoạn dây còn lại dài là:}{$12{,}75\,\text{m}$}{$13{,}75\,\text{m}$}{$12{,}65\,\text{m}$}{$13{,}65\,\text{m}$}

\vspace{0.4mm}
\noindent\textbf{\large PHẦN II. TỰ LUẬN (7,0 điểm)}

\noindent\textbf{Bài 1 (2,0 điểm). Đặt tính rồi tính:}
\vspace{0.4mm}

\noindent
\begin{tabularx}{\linewidth}{|X|X|X|X|}
\hline
\oDatTinh[4]{a)}{$58{,}74 + 26{,}49$} & 
\oDatTinh[4]{b)}{$78{,}5 - 36{,}48$} & 
\oDatTinh[4]{c)}{$162{,}38 + 54{,}7$} & 
\oDatTinh[4]{d)}{$95 - 42{,}65$} \tabularnewline
\hline
\end{tabularx}

\newpage
% -------------------- TRANG 2: MẶT SAU --------------------
\noindent
\begin{tabularx}{\linewidth}{|X|c|}
\hline
\textbf{PHẦN TỰ LUẬN (TIẾP THEO) -- PHIẾU BÀI TẬP TUẦN 10} & \textbf{MÃ ĐỀ: 506} \\
\hline
\end{tabularx}

\vspace{0.8mm}

% BÀI 2: TÌM X (4 Ý)
\khungBaiTap{%
  \textbf{Bài 2 (1,5 điểm). Tìm $x$:}\par\vspace{0.6mm}
  \noindent
  \begin{minipage}[t]{0.485\linewidth}
    \textbf{a)} $x + 32{,}4 = 81{,}2$\par
    \dongke{3}
    \vspace{0.4mm}
    \textbf{c)} $95{,}2 - x = 53{,}6$\par
    \dongke{3}
  \end{minipage}\hfill
  \begin{minipage}[t]{0.485\linewidth}
    \textbf{b)} $x - 41{,}5 = 26{,}8$\par
    \dongke{3}
    \vspace{0.4mm}
    \textbf{d)} $x + 9{,}2 = 28{,}4 - 8{,}6$\par
    \dongke{3}
  \end{minipage}
}

\vspace{0.8mm}

% BÀI 3: ĐỔI ĐƠN VỊ (4 Ý - 1 CỘT TỪ TRÊN XUỐNG)
\khungBaiTap{%
  \textbf{Bài 3 (1,0 điểm). Viết số thích hợp vào chỗ chấm:}\par\vspace{0.8mm}
  \noindent
  a) $9\,\text{m}\,52\,\text{cm} = \dotfill\,\text{m}$\par\vspace{1.4mm}
  \noindent
  b) $7\,\text{km}\,480\,\text{m} = \dotfill\,\text{km}$\par\vspace{1.4mm}
  \noindent
  c) $8\,\text{kg}\,125\,\text{g} = \dotfill\,\text{kg}$\par\vspace{1.4mm}
  \noindent
  d) $5\,\text{tấn}\,65\,\text{kg} = \dotfill\,\text{tấn}$\par
}

\vspace{0.8mm}

% BÀI 4: BÀI TOÁN GIẢI (11 DÒNG KẺ)
\khungBaiTap{%
  \textbf{Bài 4 (1,5 điểm).} Một xe tải chở xi măng. Chuyến đầu xe chở 3{,}8 tấn, chuyến thứ hai chở nhiều hơn chuyến đầu 0{,}65 tấn. Sau hai chuyến, xe chở thêm 2{,}4 tấn nữa. Hỏi xe đã chở tất cả bao nhiêu tấn xi măng?
  \par\vspace{0.3mm}%
  \noindent\textit{Bài giải:}
  \dongke{11}
}

\vspace{0.8mm}

% BÀI 5: TOÁN NÂNG CAO (8 DÒNG KẺ)
\khungBaiTap{%
  \textbf{Bài 5 (1,0 điểm).} Tính bằng cách thuận tiện nhất:\par\vspace{0.6mm}\centerline{$F=18{,}6-4{,}25-5{,}75$}
  \par\vspace{0.3mm}%
  \noindent\textit{Bài làm:}
  \dongke{8}
}

\newpage

% =============================================================================
% ========================= MÃ ĐỀ: 507 ========================================
% =============================================================================

% -------------------- TRANG 1: MẶT TRƯỚC --------------------
\noindent
\begin{tabularx}{\linewidth}{|X|c|}
\hline
\textbf{PHIẾU BÀI TẬP CUỐI TUẦN -- TOÁN LỚP 5 (TUẦN 10)} & \textbf{MÃ ĐỀ: 507} \\
\textit{Trung tâm Toán Thầy Lĩnh -- Năm học: 2026 - 2027} & \textit{Thời gian: 45 phút} \\
\hline
\end{tabularx}

\vspace{1mm}

\noindent
\begin{tabularx}{\linewidth}{|X|p{3.5cm}|}
\hline
\rule[-2mm]{0pt}{7.5mm}\textbf{Họ và tên:} \dotfill & \textbf{Lớp:} \makebox[2cm]{\dotfill} \\
\hline
\end{tabularx}

\vspace{1mm}

\noindent
\begin{tabularx}{\linewidth}{|p{3cm}|X|}
\hline
\centering\textbf{ĐIỂM SỐ} & \centering\textbf{LỜI NHẬN XÉT CỦA THẦY CÔ} \tabularnewline
\hline
\rule{0pt}{1.1cm} & \tabularnewline
\hline
\end{tabularx}

\vspace{1mm}

\noindent\textbf{BẢNG ĐIỀN ĐÁP ÁN TRẮC NGHIỆM (Mỗi câu đúng 0,25 điểm):}
\vspace{0.6mm}

\begin{tabularx}{\linewidth}{|c|*{12}{>{\centering\arraybackslash}X|}}
\hline
\textbf{Câu} & \textbf{1} & \textbf{2} & \textbf{3} & \textbf{4} & \textbf{5} & \textbf{6} & \textbf{7} & \textbf{8} & \textbf{9} & \textbf{10} & \textbf{11} & \textbf{12} \\
\hline
\textbf{Đáp án} & & & & & & & & & & & & \\
\hline
\end{tabularx}

\vspace{1.2mm}

\noindent\textbf{\large PHẦN I. TRẮC NGHIỆM KHÁCH QUAN (3,0 điểm)}

\cauTN{1}{Kết quả của phép tính $54{,}86 + 37{,}49$ là:}{$92{,}35$}{$91{,}35$}{$92{,}25$}{$93{,}35$}
\cauTN{2}{Kết quả của phép tính $96{,}83 - 48{,}52$ là:}{$48{,}21$}{$48{,}31$}{$48{,}41$}{$47{,}31$}
\cauTN{3}{Kết quả của phép tính $88{,}3 - 47{,}56$ là:}{$40{,}84$}{$40{,}64$}{$40{,}74$}{$39{,}74$}
\cauTN{4}{Kết quả của phép tính $110 - 53{,}75$ là:}{$56{,}25$}{$57{,}25$}{$56{,}15$}{$57{,}15$}
\cauTN{5}{Biểu thức nào sau đây bằng $12{,}5 + 8{,}7$:}{$8{,}7 - 12{,}5$}{$8{,}7 + 12{,}5$}{$12{,}5 - 8{,}7$}{$12{,}5 \times 8{,}7$}
\cauTN{6}{Giá trị của biểu thức $41{,}5 + 23{,}7 + 8{,}5$ là:}{$72{,}7$}{$73{,}2$}{$73{,}7$}{$74{,}7$}
\cauTN{7}{Tính nhanh tổng $9{,}1 + 4{,}35 + 0{,}9 + 5{,}65$ được kết quả là:}{$20$}{$21$}{$19{,}5$}{$20{,}9$}
\cauTN{8}{Tìm $x$, biết: $x + 42{,}5 = 96{,}8$. Giá trị của $x$ là:}{$53{,}3$}{$54{,}3$}{$55{,}3$}{$54{,}5$}
\cauTN{9}{Tìm $x$, biết: $105{,}4 - x = 62{,}8$. Giá trị của $x$ là:}{$41{,}6$}{$43{,}2$}{$42{,}6$}{$42{,}4$}
\cauTN{10}{Điền dấu thích hợp: $56{,}4 + 18{,}6 \dots 92{,}5 - 17{,}2$:}{$<$}{$>$}{$=$}{Không so sánh được}
\cauTNhaiDong{11}{Một đội công nhân làm đường trong hai ngày. Ngày thứ nhất sửa được $145{,}5\,\text{m}$, ngày thứ hai sửa nhiều hơn ngày thứ nhất $28{,}6\,\text{m}$. Cả hai ngày đội sửa được là:}{$318{,}6\,\text{m}$}{$319{,}6\,\text{m}$}{$320{,}6\,\text{m}$}{$315{,}5\,\text{m}$}
\cauTNhaiDong{12}{Chu vi hình tam giác có độ dài các cạnh lần lượt là $14{,}2\,\text{cm};\; 18{,}6\,\text{cm};\; 15{,}4\,\text{cm}$ là:}{$48{,}2\,\text{cm}$}{$47{,}2\,\text{cm}$}{$49{,}2\,\text{cm}$}{$48{,}8\,\text{cm}$}

\vspace{0.4mm}
\noindent\textbf{\large PHẦN II. TỰ LUẬN (7,0 điểm)}

\noindent\textbf{Bài 1 (2,0 điểm). Đặt tính rồi tính:}
\vspace{0.4mm}

\noindent
\begin{tabularx}{\linewidth}{|X|X|X|X|}
\hline
\oDatTinh[4]{a)}{$63{,}58 + 29{,}67$} & 
\oDatTinh[4]{b)}{$88{,}3 - 47{,}56$} & 
\oDatTinh[4]{c)}{$175{,}46 + 48{,}9$} & 
\oDatTinh[4]{d)}{$110 - 53{,}75$} \tabularnewline
\hline
\end{tabularx}

\newpage
% -------------------- TRANG 2: MẶT SAU --------------------
\noindent
\begin{tabularx}{\linewidth}{|X|c|}
\hline
\textbf{PHẦN TỰ LUẬN (TIẾP THEO) -- PHIẾU BÀI TẬP TUẦN 10} & \textbf{MÃ ĐỀ: 507} \\
\hline
\end{tabularx}

\vspace{0.8mm}

% BÀI 2: TÌM X (4 Ý)
\khungBaiTap{%
  \textbf{Bài 2 (1,5 điểm). Tìm $x$:}\par\vspace{0.6mm}
  \noindent
  \begin{minipage}[t]{0.485\linewidth}
    \textbf{a)} $x + 36{,}5 = 92{,}4$\par
    \dongke{3}
    \vspace{0.4mm}
    \textbf{c)} $102{,}5 - x = 58{,}4$\par
    \dongke{3}
  \end{minipage}\hfill
  \begin{minipage}[t]{0.485\linewidth}
    \textbf{b)} $x - 45{,}8 = 32{,}7$\par
    \dongke{3}
    \vspace{0.4mm}
    \textbf{d)} $x + 11{,}5 = 32{,}6 - 9{,}4$\par
    \dongke{3}
  \end{minipage}
}

\vspace{0.8mm}

% BÀI 3: ĐỔI ĐƠN VỊ (4 Ý - 1 CỘT TỪ TRÊN XUỐNG)
\khungBaiTap{%
  \textbf{Bài 3 (1,0 điểm). Viết số thích hợp vào chỗ chấm:}\par\vspace{0.8mm}
  \noindent
  a) $12\,\text{m}\,85\,\text{cm} = \dotfill\,\text{m}$\par\vspace{1.4mm}
  \noindent
  b) $9\,\text{km}\,650\,\text{m} = \dotfill\,\text{km}$\par\vspace{1.4mm}
  \noindent
  c) $5\,\text{kg}\,375\,\text{g} = \dotfill\,\text{kg}$\par\vspace{1.4mm}
  \noindent
  d) $8\,\text{tấn}\,90\,\text{kg} = \dotfill\,\text{tấn}$\par
}

\vspace{0.8mm}

% BÀI 4: BÀI TOÁN GIẢI (11 DÒNG KẺ)
\khungBaiTap{%
  \textbf{Bài 4 (1,5 điểm).} Một đoạn đường dài 8{,}5 km. Đội Một làm được 2{,}75 km, đội Hai làm được nhiều hơn đội Một 0{,}5 km. Hỏi đội Ba còn phải làm bao nhiêu ki-lô-mét nữa để hoàn thành đoạn đường?
  \par\vspace{0.3mm}%
  \noindent\textit{Bài giải:}
  \dongke{11}
}

\vspace{0.8mm}

% BÀI 5: TOÁN NÂNG CAO (8 DÒNG KẺ)
\khungBaiTap{%
  \textbf{Bài 5 (1,0 điểm).} Tính bằng cách thuận tiện nhất:\par\vspace{0.6mm}\centerline{$G=24{,}5-6{,}8-3{,}2$}
  \par\vspace{0.3mm}%
  \noindent\textit{Bài làm:}
  \dongke{8}
}

\newpage
\end{document}
